\documentclass[onecolumn, draftclsnofoot,10pt, compsoc]{IEEEtran}
\usepackage{graphicx}
\usepackage{url}
\usepackage{setspace}

\usepackage{geometry}
\geometry{textheight=9.5in, textwidth=7in}

\usepackage{pgfgantt}
\definecolor{grey}{rgb}{0.57, 0.64, 0.69}

% 1. Fill in these details
\def \CapstoneTeamName{The Cleverly Named Team}
\def \CapstoneTeamNumber{20}
\def \GroupMemberOne{Cade Raichart}
\def \GroupMemberTwo{Daniel Kato}
\def \GroupMemberThree{Nathan Shepherd}
\def \CapstoneProjectName{ISS Barometer App}
\def \CapstoneSponsorCompany{NASA}
\def \CapstoneSponsorPerson{Don Pettit}

% 2. Uncomment the appropriate line below so that the document type works
\def \DocType{	%Problem Statement
				% Requirements Document
				%Technology Review
				Design Document
				%Progress Report
				}

\newcommand{\NameSigPair}[1]{\par
\makebox[2.75in][r]{#1} \hfil 	\makebox[3.25in]{\makebox[2.25in]{\hrulefill} \hfill		\makebox[.75in]{\hrulefill}}
\par\vspace{-12pt} \textit{\tiny\noindent
\makebox[2.75in]{} \hfil		\makebox[3.25in]{\makebox[2.25in][r]{Signature} \hfill	\makebox[.75in][r]{Date}}}}
% 3. If the document is not to be signed, uncomment the RENEWcommand below
%\renewcommand{\NameSigPair}[1]{#1}

%%%%%%%%%%%%%%%%%%%%%%%%%%%%%%%%%%%%%%%
\begin{document}
\begin{titlepage}
    \pagenumbering{gobble}
    \begin{singlespace}
    	\includegraphics[height=4cm]{coe_v_spot1}
        \hfill
        % 4. If you have a logo, use this includegraphics command to put it on the coversheet.
        %\includegraphics[height=4cm]{NASA}
        \par\vspace{.2in}
        \centering
        \scshape{
            \huge CS Capstone \DocType \par
            {\large\today}\par
            \vspace{.5in}
            \textbf{\Huge\CapstoneProjectName}\par
            \vfill
            {\large Prepared for}\par
            \Huge \CapstoneSponsorCompany\par
            \vspace{5pt}
            {\Large\NameSigPair{\CapstoneSponsorPerson}\par}
            {\large Prepared by }\par
            Group\CapstoneTeamNumber\par
            % 5. comment out the line below this one if you do not wish to name your team
            %\CapstoneTeamName\par
            \vspace{5pt}
            {\Large
                \NameSigPair{\GroupMemberOne}\par
                \NameSigPair{\GroupMemberTwo}\par
                \NameSigPair{\GroupMemberThree}\par
            }
            \vspace{20pt}
        }
        \begin{abstract}
        % 6. Fill in your abstract
	This document focuses on the individual components of our project and “how” they will be implemented. 
	We chose to focus on the design viewpoints in terms of the using the IEEE std 1016-2009 document for software design description, we have chosen the 6 view points we chose as suitable for our project
        \end{abstract}
    \end{singlespace}
\end{titlepage}
\newpage
\pagenumbering{arabic}
\tableofcontents
% 7. uncomment this (if applicable). Consider adding a page break.
%\listoffigures
%\listoftables
\clearpage

\section{Introduction}
\subsection{Purpose}
The \CapstoneProjectName will provide crew members of the International Space Station with a readout of barometric pressure, as read from the iPads internal sensor.
This application will supplement the current method of air pressure reading, the MANOVACUMETER, a mechanical barometer.
As our application is targeted toward the crew members, it needs to be highly readable, clear, and good at quickly conveying data.
Along with ease of use, the need for readability and efficiency of use are the main decision factors in our production, as well as our design \cite{probStat}.

\subsection{Scope}
Listed below are the final product requirements of this product; version one will be provided at the end of December 2017, version two will be provided before the end of May.
The optional research, referred to as O1, will be completed if time allows, and is not considered necessary for completion of the product.

\begin{itemize}
\item[V1:] Data Screen with initial pressure, current pressure, start/stop button, time stamp, pressure change (measured in sec / mmHg)
\item[V1:] Graph Screen with plot pressure as a function of time
\item[V1:] Additional settings menu that includes plot scale, scroll functionality, orientation, decimal point, as well as other other units of measure for pressure
\item[V1:] All items will be optimized for efficient reading and clear interface
\item[V2:] Graph will have pinch to zoom functionality, scaling options (to be set in settings)
\item[V2:] Graph will be able to be exported as a csv, excel readable file
\item[V2:] Graph can be compared to t-Res table, to measure time remaining
\item[O1:] Research can be done to identify low pressure accuracy and resistance
\end{itemize}

As approved by our client, these are the goals of our project, with unlisted items being out of scope.
This may be changed in the future, depending on available time, for instance, the research portion may not be feasible, due to time restrictions.
Our version one application will be given to the client, who will then be able to give it to the ISS crew members.
Although they are the intended target user, giving them version one will allow our team to improve upon any issues that the find in using the application.
This means that our version two application may have changes not mentioned in the list, as they may be specific to the users and will only be apparent after use.

\subsection{Context}
This application is being developed for use on the ISS.
The stakeholders of the \CapstoneProjectName include the client, Donald Pettit, and the current and future crew members of the ISS, as well as any NASA employees that will benefit from this application.
That includes Ground Control members, as well as testers that may authorize our application for installation.
Although there are other applications that read and record air pressure data, they have been deemed unfit for the following reasons.
\begin{itemize}
	\item Complex and hard-to-read UI
	\item Imprecise readings or recordings
	\item Lack of customizability involving units
	\item Optimized for altitude or weather
\end{itemize}
These needs prompted this project, and serve as simplified concerns that we seek to avoid.
The users of the application, crew members of the ISS, require that the application meet their needs, as listed in the Scope.
Because we will be avoiding pitfalls of other invalid applications, our application will be tailored specifically to serve the crew members.
This means that their feedback, after version one is supplied, will be invaluable in helping guide changes to be made in version two.
As a result, we should produce an application that is designed from the bottom up, for the crew members and NASA to use in the future.

\section{Timeline}
The following is a chart that displays our projected timeline \cite{reqDoc}.
\vspace{20pt} \\
\begin{ganttchart}[
	hgrid,
    x unit=.5mm,
    time slot format=isodate,
    link bulge=4
]{2017-09-10}{2018-06-20}
% \gantttitle{\CapstoneProjectName}
\gantttitlecalendar{year, month} \\

%Items
\ganttgroup{Version 1}{2017-09-20}{2017-12-30} \\ %Will add new line
\ganttbar{Display Readings}{2017-10-30}{2017-11-20} \\
\ganttbar{Basic Plot}{2017-11-20}{2017-11-30} \\
\ganttbar{Basic Settings}{2017-11-20}{2017-11-30} \\
\ganttbar{Finalize Design}{2017-12-01}{2018-01-01} \\
\ganttgroup{Version 2}{2018-01-01}{2018-05-18} \\
\ganttbar[inline]{Test Accuracy}{2018-01-01}{2018-05-18} \\
\ganttbar[inline]{Adv. Plot}{2018-01-01}{2018-03-20} \\
\ganttbar[inline]{Adv. Settings}{2018-01-01}{2018-03-10} \\
\ganttbar[inline]{T-res Calculation}{2018-01-01}{2018-03-10} \\
\ganttbar[
	bar/.append style={fill=grey, rounded corners=3pt}
    ]{Final Submission}{2018-05-18}{2018-06-10}

%Relationships
\ganttlink{elem1}{elem2}
\ganttlink[link type=s-s]{elem2}{elem3}

\ganttlink{elem4}{elem5}
\ganttlink[link type=s-s]{elem6}{elem7}
\ganttlink[link type=s-s]{elem7}{elem8}
\ganttlink[link type=s-s]{elem8}{elem9}

\ganttlink{elem5}{elem10}

\end{ganttchart}

% Super hacky but pretty glossary
\section{Glossary of Terms}
\begin{itemize}
\item[] \textbf{2.1 csv:} A file format readable by Excel and other spread sheet software
\item[] \textbf{2.2 iOS:} The internal operating system of Apple products, specifically iPads
\item[] \textbf{2.3 iPad:} A touch screen tablet device with barometer capabilities
\item[] \textbf{2.4 ISS:} International Space Station
\item[] \textbf{2.5 MANOVACUMETER:} The current mechanical air pressure gauge aboard the ISS
\item[] \textbf{2.6 mmHg:} Millimeters of Mercury, the standard pressure unit
\item[] \textbf{2.7 NASA:} National Aeronautics and Space Administration
\item[] \textbf{2.8 t-Res:} The chart that NASA uses to calculate residual time left on station
\item[] \textbf{2.9 UI:} Acronym for user interface
\end{itemize}

\section{References}
\begingroup
\renewcommand{\section}[2]{}%
%\renewcommand{\chapter}[2]{}% for other classes
\bibliographystyle{IEEEtran}
\bibliography{bibliography}
\endgroup
\section{Context Viewpoint}
% The Context view of a system defines the relationships, dependencies, and interactions between the system and its environment—the people, systems, and external entities with which it interacts.
\subsection{Design Concerns}
% Could be setting of use [in space, iss, quick pressure change], who uses it [crew members]...
The context viewpoint is a project viewpoint that looks at the application in respect to its relationship with the environment.
For this viewpoint, the major design concerns are related to the applications use, its users, and its location.
The stakeholders of this perspective are the end users, the client, and NASA as a whole.
All of them rely on the ability of our application to be clear in how it interacts through its UI.
The end users require that the \CapstoneProjectName provide accurate data clearly and quickly, with easily understood interactivity.
The users will be reliant on this application in potentially dire situations, leaks in the ISS, meaning it must not be unduly difficult to understand or operate.
As mentioned in context section, the problem of complicated UI and difficult to understand displays are the reasons that other applications are not used.
NASA requires that this application not introduce any security problems, and that any data displayed can be trusted.
NASA also requires that the application be functional and clear of critical bugs, which could jeopardize mission safety.
The latter requirement is lessened by the fact that NASA will keep the MANOVACUMETERS aboard the ISS, as well as other redundancies, but should still be considered.
The setting of use, the ISS, introduces some interesting considerations to the design.
For instance the accelerometer will not work as it does in 1g, and requires shaking to reorient the iPad.
Also the iPad may be exposed to extreme low pressures, and should be trusted to remain functional and accurate.
These considerations do not affect the design of the application, but may be tested if the optional goals are achieved.
In addition to low pressure and gravity, the stations remoteness may also affect design, in that updates and communications will be sent through NASA and their network.

\subsection{Design Elements}
% UI design, barometer implementation, Platform(iPad vs ...)
In response to the design concerns, several design elements have been chosen to alleviate the dilemmas.
As a result of the crew members need for clear, usable design, we have decided to use Swift Charts as our plot library, and a combination of Cocoa controls, FlatUIKit, and  CRPageViewController \cite{techDocRaichart}.
In addition to specific UI libraries, we plan to design the UI with simplicity and readability in mind.
This will be judged in part by our client, and eventually by the crew members once it is deployed, as version one.
The plot library will provide the functionality that is needed, in addition to being a supported and trusted plotter.
Using trusted and recommended libraries and packages will also help provide NASA with understandable and trusted code.
This means that their approval and authorization will be quicker and therefore allow faster turn around on things such as updates.
Accelerometer issues caused by low gravity is an easily solvable issue, requiring only the addition of a button to switch orientation.
This has been a design choice, proposed by the client, since the beginning, and would solve the issue for the end users.
The other problem caused by the environment, that of low pressure accuracy, will not be as easily solved.
As our application is not able to alter the hardwares capabilities in withstanding pressure, the only option would be to alter the platform.
The options that we have considered include external Arduino applications, the iPhone 6, or the iPad Air.
The iPad has been identified as our required platform, by our client, and backed by our research \cite{techDocRaichart}.
This means our only option for ensuring the reliability of the application is to test it in a controlled pressure container.
This would require things outside the scope of this project, but may be achieved as a stretch goal, as seen in the Scope section.
The added difficulty of transferring data through NASA to the ISS means that updates will be sparse, and return data or remote data storage is infeasible.
These factors guided our decision to use CoreData and NSUserDefaults as our storage packages, instead of a remote server package, like SQLite \cite{techDocShepherd}.
The choices made for the Context Viewpoint will help optimize our application for the location and users involved.

\section{Functional Viewpoint}
% The Functional view of a system defines the architectural elements that deliver the system’s functionality.
\subsection{Design Concerns}
% Could be accuracy and precision, speed, clarity of graph...
\subsection{Design Elements}
% Could be UI Design, graph interface, language, storage information, barometer implementation...

\section{Information Viewpoint}
% Describes the way that the system stores, manipulates, manages, and distributes information
\subsection{Design Concerns}
% Could mention accessibility to csv, speed, security from errors
The stakeholders of the information viewpoint are those reliant on accurate and secure data storage.
This is all of the main stakeholders: the client, the ISS crew members, our team, and NASA, as a whole.
As the purpose of our application is to display information from the barometer module of the iPad, so it follows that it is necessary to display correct and timely data.
Also necessary, is the ability to store the data between uses, so that a user may exit the application and open to the same view, complete with prior data.
One additional concern may be the need to record data while the application is not running in the foreground of the iPad.
This may be a concern, but from our understanding of the clients specifications, this is not within scope, and shouldn't be an issue.
Once the user has a plot recorded they might want to export it as a csv file, which must also be secure from errors or tampering.
It must be ensured that the data being exported to the csv is the data that was recorded.
Once the csv is exported, that data is no longer our concern, as it is in the hands of the user.
In addition to exporting and saving plot data between uses, our application will have some various settings, which will be stored between uses, so that certain settings are persistent.
The settings that must remain persistent between uses are listed below \cite{reqDoc}:
\begin{itemize}
	\item Significant digits
	\item Plot scale
	\item Plot scroll functionality
	\item Orientation
	\item Units of pressure
	\item Sample rate of recording
\end{itemize}

\subsection{Design Elements}
% Essentially just the graph storage function, and settings storage
The main security that will be implemented is done through using well defined and secure storage packages.
We are not expecting intentional tampering, as the application will not be connected to the network, and the only users are using the application for self-preservation.
That said, it is still possible that data be mangled, dropped, or altered as it is recorded or exported.
By correctly using well maintained and trusted packages we will avoid many problems with data storage.
As the data being stored is small, and relatively simple, the use of packages will simplify our storage, while allowing the users quick and reliable data.
The packages that will store plot data and settings data will be two different packages.
This method, as suggested by Apple\cite{coreApple}, will separate user data, such as settings, from application data, like the recordings or the plot.
The user data package will be NSUserDefaults, which is suggested for storing small portions of data.
The settings stored will be integers or boolean data, and like other application data, will be stored locally on the iPad.
The plot storage package will be CoreData, which will be able to handle larger sets of data, like the plot.
This also will be stored locally.

\subsection{Design Attributes}
Data stored by the NSUserDefaults will be of the form of integers or booleans, as it will only be storing user choices.
Data stored by CoreData, for the plot, will be in the form of paired floats and timestamps.
Together each of these data points will form a data point on the plot, as the recording continues, and the graph implementation plots the points.
If the data is recorded and then exported, the user will be prompted whether or not to remove the data stored in the application.
If the user accepts then the plot data will be removed, and the screen will be returned to its prerecord state.
The data must be removed from the application, either by deleting or exporting, before the user will be allowed to record again.
This technique prevents overflow of data storage, which would cause the application to crash.

\section{Concurrency Viewpoint}
% Describes the concurrency structure of the system, mapping functional elements to concurrency units to clearly identify the parts of the system that can execute concurrently, and shows how this is coordinated and controlled
\subsection{Design Concerns}
% Will just probably be graph running while app closed? Barometer keeping up to date, no locks on data so not really a problems
For this viewpoint the major concerns are related to the speed of hardware, and timing.
The stakeholders of this would be our client, end users and astronauts aboard the ISS. 
The astronauts will need to rely on our app to correctly have all the timing work in unison.  
When using hardware and software systems together we run into some timing problems. 
Our app might have a timing concern with how fast the barometer records atmospheric pressure. 
The barometer might be reading the data slower than the graph can display it, also it might be reading it faster than we can display the data on the graph.
This can cause problems with the end users. 
Looking at the graph with a older atmospheric pressure than the one the barometer can cause the user to not evacuate early, this could result in death of the crew. 
Our app will be doing many actions all at the same time, this all needs to be done simultaneously as to ensure the accuracy of the data.
Astronauts need to have all data be precise and completely accurate to save lives.  

\subsection{Design Elements}
% Barometer implementation, chart library
In response to the design concerns, several hardware chips are being used correctly calculate the timing problems. 
The ISS Barometer App will be using the iPad Air 2 for its platform. The Ipad uses a Apple A8x triple core 1.5 GHz Typhoon cpu. 
The triple core cpu allows for multiple action to be done simultaneously. 
This means that we can display the atmospheric pressure, dp/dt, calculate time to evacuation and display the graph simultaneously making it accurate and safe.
With each core of the cpu running at 1.5 GHz we do not have to worry about the barometer recording data faster than it can be displayed. 


\section{Development Viewpoint}
% Describes the architecture that supports the software development process
\subsection{Design Concerns}
% We're the only stakeholders in this one pretty much
\subsection{Design Elements}
% Language Used, IDE, packages in general

\section{Deployment and Post-Deployment Viewpoints}
% Pretty small and similar so I grouped them, easy to ungroup
% Deployment: Describes the environment into which the system will be deployed, including the dependencies the system has on its runtime environment
% Operation: Describes how the system will be operated, administered, and supported when it is running in its production environment
\subsection{Design Concerns}
This viewpoint is a project that looks at the environment that the system will be deployed in. 
The operational concerns are how  it will be operated, administered and supported while in deployment.
The second concern of this viewpoint is will the app help the crew while in use. 
The ISS Barometer App will be aboard the ISS space station. 
The space station is orbiting the Earth at a minimum of 330 km. 
The interior of the ISS has a pressure of 1 atmosphere, measured in millimeters of Mercury (mmHg), that must be maintained for the safety of the Astronauts aboard the ISS.
If the atmospheric pressure drops below .540 atm it will endangering the crew. 
Leaks in the ISS can be caused by faulty seals in hatches or orbital debris striking the vessel, a severe leak consists of a loss of 1 mm per minute.
\subsection{Design Elements}
In response to the viewpoint concerns our app will be deployed to find these leaks. 
The astronauts will operate the app by going from room to room, closing hatch door behind them and seeing if the atmospheric pressure stabilizes. If it doesn't stabilize they will move on to the next room. 
Until they find the room/rooms with the leak or they must abandon ship. When the app is activated with the start button, it will record the atmospheric pressure and display it along with the appropriate time stamp. 
That data will be used to calculate the dp/dt. dp/dt will be displayed by the first page, used as the slope of the graph, and it will be used with the comparison to the t-Res tables to display the time till evacuation. 
The App will be administered by uploading the files to the space station and downloading onto the iPads.
Updates to the app will administered to the iPad in the same way. As the iPad's OS gets update so must we update the app.
\section{Conclusion}
As described in our Technology documents, we have sought out different options for the major pieces of this project.
With the research done, and the viewpoints provided here, it has been shown that the \CapstoneProjectName is properly thought out and detailed.
This application is simple in its infrastructure, but has the potential for quality design and efficient and reliable displays.
In addition, testing this application, and the limitations of the iPad, if completed, may provide beneficial insight to the operating limitations and reliability of the readings.
These improvements are planned for version two, as mentioned in the Scope.
Using the viewpoints from this document, in addition to potential feedback from the users after version one, there may be many unforeseen improvements to be made.
As those improvements are made, our application will become perfectly tailored to the users' needs.


\end{document}
